%%%%%%%%%%%%%%%%%%%%%%%%%%%%%%%%%%%%%%%%%%%%%%%%%%%%%%%%%%%%%%%%%%%%%%
% Lớp 10 | Phương pháp tọa độ trong mặt phẳng (𝑂𝑥𝑦) | Đường thẳng trong mặt phẳng tọa độ
% Dạng 0H9?3-8: Bài toán thực tế, PP tọa độ hóa
% Số câu: 5  (sinh tự động bởi tach_cau.py; từ nay thầy sửa trực tiếp file này)
%%%%%%%%%%%%%%%%%%%%%%%%%%%%%%%%%%%%%%%%%%%%%%%%%%%%%%%%%%%%%%%%%%%%%%

%%%%%%%%%%%%%%%% ĐÚNG/SAI %%%%%%%%%%%%%%%%

% nguồn: MaHoaNH_DungSai_Lop10
\begin{ex}%[0H9N3-8]%[ma:76b4837d]
	Chuyển động của vật thể $M$ được thể hiện trên mặt phẳng toạ độ $Oxy$. Vật thể $M$ khởi hành từ điểm $A(5; 3)$ và  chuyển động thẳng đều với vectơ vận tốc là $\overrightarrow{v}(1; 2)$. Các mệnh đề sau đúng hay sai?
	\def\dotEX{} 
	 \choiceTF
		{\True Vectơ chỉ  phương của đường thẳng biểu diễn chuyển động của vật thể là $\overrightarrow{v}(1; 2)$.}
		{\True Vật thể $M$ chuyển động trên đường  thẳng $2x-3y-1=0$.}
		{Toạ độ của vật thể $M$ tại thời điểm $t$ $(t > 0)$ tính từ  khi khởi hành là $\heva{&x=5+t \\&y=3+2t.}$}
		{\True Khi $t=5$ thì vật thể  $M$ chuyển động được quãng đường dài bằng $5\sqrt{5}$}
	 \loigiai{
		\begin{itemchoice}
			\itemch vectơ chỉ phương của đường thẳng  biểu diễn chuyển động của vật thể là $\overrightarrow{v}(1; 2)$, do đó đường thẳng này có vectơ pháp tuyến là $\overrightarrow{n}(2;-1)$.
			\itemch Mặt  khác, đường thẳng
			này đi qua điểm $A(5; 3)$ nên có phương  trình là
			 \[2(x-5)-(y-3)=0\Leftrightarrow 2x-y-7=0.\]
			\itemch Vật thể khởi hành từ điểm $A(5; 3)$ và chuyển động thẳng đều với vectơ  vận tốc
			là $\overrightarrow{v}=(1; 2)$ nên vị trí của vật thể tại thời điểm $t$ $(t > 0)$ có  toạ độ là $\heva{&x=5+t \\&y=3+2t.}$
			\itemch Gọi $B$ là vị trí của  vật thể tại thời điểm $t=5$. Do đó, toạ độ của điểm $B$ là
			 $$\heva{&x_B=5+5=10\\&y_B=3+2\cdot 5=13.}$$
			Khi đó quãng đường  vật thể đi được là $AB=\sqrt{25+100}=5\sqrt{5}$.
		\end{itemchoice}
	} 
\end{ex}


%%%%%%%%%%%%%%%% TRẢ LỜI NGẮN %%%%%%%%%%%%%%%%

% nguồn: MaHoaNH_TraLoiNgan_Lop10
\begin{bt}%[0H9V3-8]%[ma:cc2370a7]
	Một con tàu muốn xuất phát từ hòn đảo $A$ trở về bờ biển sau đó di chuyển đến hòn đảo $B$. Trên màn hình ra đa của trạm điều khiển  (được coi như mặt phẳng $Oxy$), vị trí điểm $A$, $B$ có tọa độ lần lượt là $A(0;0)$, $B(5;-1)$, giả sử đường bờ biển có phương trình đường thẳng là $\Delta\colon x-y+3=0$. Tọa độ điểm $M(a;b)$ trên bờ biển mà tàu sẽ di chuyển đến sao cho độ dài đường đi của tàu từ $A$ đến $B$ là ngắn nhất. Xác định tổng $S=a+b$.
	\shortans{1} 
	\loigiai{ 
		Ta nhận thấy hai điểm $A$, $B$ nằm về cùng một phía của đường  thẳng $\Delta\colon x-y+3=0$.\\
		Gọi $A'$  là điểm đối xứng của $A$ qua $\Delta$.
		\begin{center}
			\begin{tikzpicture}[line join=round,line cap=round,>=stealth,scale=0.75,font=\footnotesize]
					\foreach \x/\y/\n in {0/0/H,2/0/M,6/0/d,5/3/B,0/2/A,0/-2/A'} \coordinate (\n)  at (\x,\y);
					\draw  (A)--(A')--(B)
						(d) node[above left]{$\Delta$}--($(H)-(0.5,0)$); 
					\draw  (d) node[below left]{bờ biển};
					\draw (A)--(H) node[midway,left]{$d$}; 
					 \foreach \t/\g in {A/180,B/-90,A'/0,H/-140,M/120}{
						\fill (\t) circle (1pt)  node[shift={(\g:8pt)}]{$ \t $};
					 }
					\draw pic[thin,draw,angle radius=2mm]  {right angle= d--H--A};
				\end{tikzpicture}
		\end{center}
		Gọi $d$ là đường thẳng đi qua $A$ và vuông góc với $\Delta$ tại  $H$.\\
		Phương trình tham số  của $d$ là $\heva{&x=t\\&y=-t.}$\\
		Vì $H\in d$ nên $H(x_H;-x_H)$.\\ 
		 Mặt khác, $H\in\Delta\Rightarrow x_H-(-x_H)+3=0\Leftrightarrow x_H=-\dfrac{3}{2}$. Suy ra $H\left(-\dfrac{3}{2};\dfrac{3}{2}\right)$.\\
		Vì $H$ là  trung điểm của $AA'$ nên $A'(-3;3)$.\\
		Vì $A$, $B$ cố định nên  độ dài đường đi của tàu ngắn nhất $\Leftrightarrow AM+MB$ ngắn nhất.\\
		 Ta có $AM+MB=A'M+MB\ge A'B$.\\
		Vậy $AM+MB$ ngắn nhất $\Leftrightarrow A'$, $M$, $B$ thẳng  hàng $\Leftrightarrow A'B$ cắt $\Delta$ tại $M$.\\
		Phương trình đường thẳng $A'B$ là  $x+2y-3=0$.\\
		Tọa độ điểm $M$ là nghiệm của hệ $\heva{&x-y+3=0\\&x+2y-3=0}\Leftrightarrow\heva{&x=-1\\&y=2.}$\\ 
		 Vậy $M(-1;2)$.
	 }
\end{bt}

% nguồn: MaHoaNH_TraLoiNgan_Lop10
\begin{bt}%[0H9V3-8]%[ma:d21310f1]
	Mảnh đất hình tam giác được coi là xấu trong phong thủy. Nhưng hiện nay do sự gia tăng dân số và đô thị hóa, các mảnh đất hình tam giác ngày càng nhiều và cần phải xây nhà trên các mảnh đất đó. Để hóa giải điềm xấu khi xây nhà trên các mảnh đất hình tam giác, các nhà phong thủy đã gợi ý một cách đơn giản là đặt một chiếc đèn sáng tại tâm của ngôi nhà để có thể chiếu sáng  đến tất cả các góc của ngôi nhà. Nhà bạn Hoa đang muốn xây nhà trên mảnh đất hình tam giác của mình với các kích thước là $15$ m; $18{,}4$ m; $23$ m (như hình 1) và bạn Hoa đã gắn hình dạng mảnh đất đó lên hệ trục tọa độ $Oxy$ như hình $3$. Giả sử bạn Hoa sẽ đặt chiếc đèn phong thủy ở vị trí có tọa độ $(a, b)$. Tính $T=2a-16b$.
	\begin{center}
		\begin{tikzpicture}[line join = round,line cap = round,>=stealth,font=\footnotesize,scale=1]
				\def\a{4.6} 
				\def\b{3} 
				 \def\c{3.68}
				\path (0:0)  coordinate (B)
					++(0:\a)  coordinate (C);
				 \path[name path=c1] (B) circle (\c);
				\path[name path=c2] (C)  circle (\b);
				\path[name  intersections={of=c1 and c2,by={A,D}}];
				\pgfresetboundingbox 
				\draw  [thick](A)--(B)--(C)--cycle;
				\def\f  (#1){\a*(#1)+\b}
				\def\xmin{7} 
				 \def\xmax{12}
				\def\ymin{0} 
				 \def\ymax{3}
				 \draw[->,thick] (\xmin-1,0)--(\xmax+1,0) node[below] {$x$};
				\draw[->,thick] (7,\ymin-1)--(7,\ymax+1) node[left] {$y$}; 
				\draw (7,0) node  [below left] {$O$};
				 \draw [thick] (7,0)--(9.79,2.4)--(11.6,0);
				\draw  (9.79,0.1)--(9.79,-0.1) node [below] {$14$};
				\draw (7.1,2.4)--(6.9,2.4) node  [left] {$12$};
				\draw[dashed]  (9.79,0)--(9.79,2.4)--(7,2.4);
				\node [below] at  (11.6,0) {$23$};
				\node [above]  at (9.79,2.4) {$A$};
				\node [above right]  at (11.6,0) {$B$};
				\node [below] at (2.3,-0.5) {Hình 1}; 
				\node [below]  at (9.2,-0.5) {Hình 3};
			\end{tikzpicture}
	\end{center}
	 \shortans{11}
	\loigiai{ 
		Tâm của tam giác  là tâm đường tròn ngoại tiếp tam giác, tức là giao ba đường trung trực.\\
		Gọi $M$, $N$ lần lượt là trung điểm của các cạnh  $OA$, $OB$. Suy ra $M(7; 6)$; $N\left(\dfrac{23}{2}; 0\right)$.\\
		Đường trung trực cạnh $OA$ nhận $\overrightarrow{OA}=(14; 12)$ hay $\overrightarrow{n}=(7; 6)$ làm vectơ pháp tuyến và đi qua điểm $M(7; 6)$ nên nó có  phương trình là $7(x-7)+6(y-6)=0\Leftrightarrow 7x+6y-85=0$.\\
		Đường trung trực cạnh $OB$ có phương trình là $x=\dfrac{23}{2}$.\\ 
		Tâm của tam giác $OAB$ có tọa độ là nghiệm hệ phương  trình $\heva{&7x+6y-85=0 \\&x=\dfrac{23}{2}} \Leftrightarrow\heva{&x=\dfrac{23}{2} \\&y=\dfrac{3}{4}.}$\\
		Vậy  $\heva{&a=\dfrac{23}{2} \\
		&b=\dfrac{3}{4}} \Rightarrow T=2a-16b=2\cdot \dfrac{23}{2}-16\cdot  \dfrac{3}{4}=11$.
	} 
\end{bt}


%%%%%%%%%%%%%%%% TỰ LUẬN %%%%%%%%%%%%%%%%

% nguồn: MaHoaNH_TuLuan_Lop10
\begin{bt}%[0H9H3-8]%[ma:7830dcd8]
	Cho hệ trục tọa độ $Oxy$, trong đó mỗi đơn vị độ dài tương ứng với $1$ km. Một  chiếc xe $M$ chạy trên đường thẳng từ $A(-6;-3)$ đến $B(6;6)$ với vận tốc không đổi là $50$ km/h. Biết rằng có một ra-đa đặt tại $C(5;-1)$. Hỏi trong quá trình chiếc xe $M$ di chuyển, sau bao nhiêu phút kể từ lúc xuất phát thì khoảng cách từ $C$ đến chiếc xe $M$ là ngắn nhất và khoảng cách ấy là bao nhiêu ki-lô-mét?
	\loigiai{\immini{Ta có $\overrightarrow{AB}=(12; 9)$ nên đường thẳng đi qua $A$ và  $B$ có một vectơ pháp tuyến là $\overrightarrow{n}=(3;-4)$ nên có phương trình tổng quát là \break $3(x+6)-4(y+3)=0$ hay $3x-4y+6=0$.\\
		Khoảng cách từ điểm  $C(5,-1)$ đến đường thẳng $AB$ là
		\[ 
		d = \dfrac{|3\cdot (5) - 4(-1) + 6|}{\sqrt{3^2 + 4^2}} = \dfrac{25}{5} =  5 \text{ km}.
		\] 
		 }
		{\begin{tikzpicture}[scale=0.5,font=\footnotesize,line join=round,line cap=round,>=stealth]
		 \path
			(0,0)  coordinate (B)
			++(60:5) coordinate (C) 
			(8,0) coordinate (A); 
		\coordinate  (H) at ($(A)!(C)!(B)$);
		\draw (A)--(B) 
			(A)--(C) 
			(B)--(C)--(H); 
		\foreach  \i/\g in {A/-90,B/-90,C/90,H/-90}
		\fill[black] (\i) circle (1pt)+(\g:4mm)  node[scale=1]{$\i$};
	\end{tikzpicture}
		 }
		 Ta có $AC=\sqrt{(5+6)^2+(-1+3)^2}=\sqrt{125}$.\\
		Gọi $H$  là hình chiếu của $C$ lên $AB$. Khi đó $AH=\sqrt{AC^2-CH^2}=10$.\\
		Thời gian kể từ lúc xuất phát thì  khoảng cách từ $C$ đến chiếc xe $M$ là ngắn nhất là \break $\dfrac{10}{50}=0{,}2\text{ (h)}=12\text{ (phút)}$.\\
		Khoảng cách ngắn nhất bằng  $5$ km.
	} 
\end{bt}

% nguồn: MaHoaNH_TuLuan_Lop10
\begin{bt}%[0H9V3-8]%[ma:71d17c88]
	Trong mặt phẳng tọa độ $Oxy$, một thiết bị âm thanh được phát từ vị trí $A(5;-4)$. Người ta dự định đặt một máy thu tín hiệu trên đường thẳng có phương trình: $x-y-3=0$. Khi đặt  máy tại vị trí $M(a;b)$ sẽ nhận được tín hiệu sớm nhất. Tìm tọa độ điểm $M$.
	\loigiai{ 
		Vì $M(a; b) \in d\colon x-y-3=0$ nên $a - b - 3 = 0  \Leftrightarrow b = a - 3$.\\
		$\Rightarrow$ $M(a; a -  3)$.\\
		Ta có \(AM = \sqrt{(a - 5)^2 + (a -  3 - (-4))^2} = \sqrt{2a^2 - 8a + 26}=\sqrt{2(a - 2)^2 + 18}\geq 18\).\\
		Do đó $AM$  nhỏ nhất là $AM=18$ khi \(a - 2 = 0 \Rightarrow a = 2\).\\
		Vậy tọa độ điểm $M$ để máy thu tín hiệu sớm  nhất là $M(2; -1)$.
	 }
\end{bt}
